\begin{tikzpicture}[x=1mm,y=1mm,font=\kaishu\small,
  pgm variable/.style={circle,draw=BookBlue,fill=FigureInput!8,
    minimum size=8mm,inner sep=1pt},
  factor/.style={rectangle,draw=BookTeal,fill=FigureModel!12,
    minimum size=8mm,inner sep=1pt},
  edge/.style={draw=BookInk,semithick},
  panel/.style={align=center,text width=42mm}]
  \node[panel] at (16,14) {保留二元因子};
  \node[pgm variable] (l1) at (0,0) {$X_1$};
  \node[pgm variable] (l2) at (32,0) {$X_2$};
  \node[pgm variable] (l3) at (16,-32) {$X_3$};
  \node[factor] (f12) at (16,0) {$f_{12}$};
  \node[factor] (f13) at (4,-18) {$f_{13}$};
  \node[factor] (f23) at (28,-18) {$f_{23}$};
  \draw[edge] (l1)--(f12)--(l2);
  \draw[edge] (l1)--(f13)--(l3);
  \draw[edge] (l2)--(f23)--(l3);
  \node[panel,text=BookMuted] at (16,-45) {三个$r^2$表\\因子图有环};

  \node[panel] at (83,14) {合并为三元因子};
  \node[pgm variable] (r1) at (67,0) {$X_1$};
  \node[pgm variable] (r2) at (99,0) {$X_2$};
  \node[pgm variable] (r3) at (83,-32) {$X_3$};
  \node[factor] (f123) at (83,-13) {$f_{123}$};
  \draw[edge] (r1)--(f123);
  \draw[edge] (r2)--(f123);
  \draw[edge] (r3)--(f123);
  \node[panel,text=BookMuted] at (83,-45) {一个$r^3$表\\因子图无环};
\end{tikzpicture}
